\section{Day 8: Zero--One Laws (Oct.\ 5, 2026)}
\begin{theorem}[Panchenko \S2.8: Kolmogorov $0$--$1$ law] \label{thm:kolmogorov_0-1_law}
    Let $(X_i)_{i \geq 1}$ be a sequence of random variables. Define the \textit{tail $\sigma$-algebra},
    \[ \tau = \prod_{n \geq 1} \sigma(X_n, X_{n+1}, \dots). \]
    If the $X_i$ are independent and $A \in \tau$, then $\PP(A) \in \{0, 1\}$.
\end{theorem}

As an example, set $S_n = X_1 + \dots + X_n$, where $(X_i)_{i \geq 1}$ are real random variables. Then
\[ \left\{\limsup_{n \to \infty} S_n \in (-\infty, \infty)\right\}, \left\{\lim_{n \to \infty} \frac{S_n}{n} \text{ exists}\right\} \in \tau, \]
but $\{\limsup_{n \to \infty} S_n \geq 0\} \notin \tau$. Essentially, ``if I change finitely many variables and the event doesn't change, then it is in the tail $\sigma$-algebra''.

\begin{proof}[Proof of \ref{thm:kolmogorov_0-1_law}]
    Let $A \in \tau$. We will show that $A$ is independent of $A$, i.e., $\PP(A) = \PP(A)^2$. First, if $A \in \sigma(X_1, \dots, X_m)$ and $B \in \SB_{m+1} = (X_{m+1}, \dots)$, then $A$ and $B$ are independent.
    \begin{subproof}
        Let $\SB_{m+1,\ell} = \sigma(X_{m+1}, \dots, X_\ell)$. Then it is independent from $\sigma(X_1, \dots, X_m)$ by the grouping lemma, and so $\sigma(X_1, \dots, X_m)$ and $\bigcup_{k \geq m + 1} \SB_{m+1, k}$ are independent. We may conclude by the Dynkin $\pi$--$\lambda$ theorem since $\bigcup_{k \geq m + 1} \SB_{m+1,k}$ is a $\pi$-system.
    \end{subproof}

    Now, consider that $\sigma(X_1, \dots)$ and $\tau$ are independent. This follows because, by letting $\SA_m = \sigma(X_1, \dots, X_m)$, we see that
    \[ \sigma(X_1, \dots) = \sigma\left(\bigcup_{m \geq 1} \SA_m\right), \]
    for which $\tau = \bigcap_{m \geq 1} \SB_m$, so, per our earlier proof, it follows that $\tau$ and $\bigcup_{m \geq 1} \SA_m$ are independent. By the Dynkin $\pi$--$\lambda$ theorem again, $\tau$ and $\sigma(\bigcup_{m \geq 1} \SA_m)$ are independent. Since $A$ lies in both $\tau$ and the $\sigma$-algebra generated by unions of $\SA_m$, it follows that $\PP(A) = \PP(A)^2$.
\end{proof}

\begin{corollary}
    If $X_i$ is an independent sequence, then $\lim_{n \to \infty} S_n/n$ either diverges a.s., or a.s.\ the limit exists and is non-random (i.e., equal to some fixed constant).
\end{corollary}

Consider the ``very famous'' following example. Consider $\ZZ^d$, and, for every edge, flip a coin with probability $p$ on whether to include the edge in our set or not. For fixed $p$, what is the probability that the graph has an infinite component? Clearly, the behavior of finitely many random variables doesn't change this event, so, by \ref{thm:kolmogorov_0-1_law}, if we plot this probability as a function of $p$, it will be a monotone, right-continuous function in $\ZZ^2$. For enough $d$, a result from Hara--Slade applies. It was... apparently solved by AI a month ago with specialized human input.

\begin{theorem}[Panchenko \S2.9: Hewitt--Savage $0$--$1$ law] \label{thm:hewitt-savage_0-1_law}
    Let $(X_n)_{n \geq 1}$ be a sequence of random variables. A finite permutation of $\ZZ$ is a bijection $\sigma \colon \ZZ \to \ZZ$ such that $\sigma(i) = i$ for all large enough $i$. For a sequence $X = (X_1, X_2, \dots)$, write $\sigma(X) = (X_{\sigma(1)}, X_{\sigma(2)}, \dots)$; the \textit{exchangeable $\sigma$-algebra} for $X$ is
    \[ \xi = \{A \mid X \in A \iff \sigma(X) \in A \text{ for all finite permutations $\sigma$}\}. \]
    If the $X_i$ are i.i.d., then $A \in \xi$ implies that $\PP(A)$ is $0$ or $1$.
\end{theorem}
Note that $\xi \supsetneq \tau$. As an example to demonstrate strict containment, observe $\{\sum_{i=1}^\infty X_i = 5\} \in \xi \setminus \tau$.

\newpage
The following lemma is from Durrett's appendix, but I don't know where.
\begin{lemma}
    If $\SA$ is an algebra, and $B \in \sigma(\SA)$, then for any $\eps > 0$, we can find $A_\eps \in \SA$ such that $\PP(B \Delta A_\eps) < \eps$. Note that $\Delta$ denotes symmetric difference.
\end{lemma}

We may start to prove the theorem as follows. Let $\SA = \bigcup_{n \geq 1} \sigma(X_1, \dots, X_n)$, and let $A \in \xi$. Then, for all $\eps > 0$, there exists $n \in \NN$ such that $A_n \in \sigma(X_1, \dots, X_n)$ and $\PP(A \,\Delta\, A_n) < \eps$, since $\SA$ is an algebra generated by $\sigma(X_1, \dots)$. We can write $A_n = \{(x_1, \dots, x_n) \in B_n\}$ for some $B_n \in \SB(\RR^n)$. Without loss of generality, the $X_i$ are real-valued. From here, let $A_n' = \{(X_{n+1}, \dots, X_{2n}) \in B_n\}$.

We may now claim that $\PP(A \,\Delta\, A_n) = \PP(A \,\Delta\, A_n')$. This will be done next lecture.